\documentclass[11pt,a4paper]{article}

% --- PACKAGES ---
\usepackage[utf8]{inputenc}
\usepackage[T1]{fontenc}
\usepackage{lmodern}
\usepackage[a4paper,margin=2.3cm]{geometry}
\usepackage{amsmath,amssymb,amsthm}
\usepackage{mathtools}
\usepackage{graphicx}
\usepackage{xcolor}
\usepackage[most]{tcolorbox}
\usepackage{enumitem}
\usepackage{booktabs}
\usepackage{array}
\usepackage{tikz}
\usepackage{hyperref}
\usepackage{titlesec}
\usepackage{fancyhdr}
\usepackage{parskip}

% --- COLORS (Leiden inspired) ---
\definecolor{leidenblue}{RGB}{0, 30, 90}
\definecolor{leidenlight}{RGB}{135, 152, 195}
\definecolor{accentgreen}{RGB}{180, 198, 80}
\definecolor{boxgray}{RGB}{245, 245, 250}

% --- HYPERREF ---
\hypersetup{
    colorlinks=true,
    linkcolor=leidenblue,
    urlcolor=leidenblue,
    citecolor=leidenblue,
    pdftitle={Neural Computing Study Guide},
    pdfauthor={Alessandro Garcia}
}

% --- TITLE FORMATTING ---
\titleformat{\section}
    {\color{leidenblue}\Large\bfseries\raggedright}{\thesection}{1em}{}[\titlerule]
\titleformat{\subsection}
    {\color{leidenblue}\large\bfseries\raggedright}{\thesubsection}{1em}{}
\titleformat{\subsubsection}
    {\color{leidenblue}\normalsize\bfseries\raggedright}{\thesubsubsection}{1em}{}

% --- HEADER/FOOTER ---
\pagestyle{fancy}
\fancyhf{}
\fancyhead[L]{\color{leidenblue}\small\textit{Neural Computing -- Study Guide}}
\fancyhead[R]{\color{leidenblue}\small\thepage}
\renewcommand{\headrulewidth}{0.4pt}
\renewcommand{\headrule}{\hbox to\headwidth{\color{leidenblue}\leaders\hrule height \headrulewidth\hfill}}

% --- CUSTOM BOXES ---
\newtcolorbox{definitionbox}[1][]{
    enhanced,
    colback=boxgray,
    colframe=leidenblue,
    fonttitle=\bfseries\color{white},
    coltitle=white,
    title=Definition,
    boxrule=0.5pt,
    arc=2mm,
    left=4mm,
    right=4mm,
    top=2mm,
    bottom=2mm,
    #1
}

\newtcolorbox{keybox}[1][]{
    enhanced,
    colback=yellow!8,
    colframe=orange!70!black,
    fonttitle=\bfseries\color{white},
    coltitle=white,
    title=Key Idea,
    boxrule=0.5pt,
    arc=2mm,
    left=4mm,
    right=4mm,
    top=2mm,
    bottom=2mm,
    #1
}

\newtcolorbox{theorembox}[1][]{
    enhanced,
    colback=blue!4,
    colframe=leidenblue,
    fonttitle=\bfseries\color{white},
    coltitle=white,
    title=Theorem,
    boxrule=0.5pt,
    arc=2mm,
    left=4mm,
    right=4mm,
    top=2mm,
    bottom=2mm,
    #1
}

\newtcolorbox{examplebox}[1][]{
    enhanced,
    colback=accentgreen!10,
    colframe=accentgreen!80!black,
    fonttitle=\bfseries\color{white},
    coltitle=white,
    title=Example,
    boxrule=0.5pt,
    arc=2mm,
    left=4mm,
    right=4mm,
    top=2mm,
    bottom=2mm,
    #1
}

\newtcolorbox{warningbox}[1][]{
    enhanced,
    colback=red!5,
    colframe=red!70!black,
    fonttitle=\bfseries\color{white},
    coltitle=white,
    title=Caveat / Pitfall,
    boxrule=0.5pt,
    arc=2mm,
    left=4mm,
    right=4mm,
    top=2mm,
    bottom=2mm,
    #1
}

% --- MATH SHORTCUTS ---
\newcommand{\R}{\mathbb{R}}
\newcommand{\E}{\mathbb{E}}
\newcommand{\vw}{\vec{w}}
\newcommand{\vx}{\vec{x}}
\newcommand{\vy}{\vec{y}}
\newcommand{\vh}{\vec{h}}
\newcommand{\vC}{\vec{C}}
\newcommand{\sgn}{\text{sgn}}

% --- TITLE PAGE ---
\begin{document}

\thispagestyle{empty}

\begin{center}
\vspace*{0.5cm}
{\color{leidenblue}\Huge\bfseries Neural Computing}\\[0.5em]
{\color{leidenlight}\Large Comprehensive Study Guide}\\[0.3em]
{\large \textit{Leiden University -- LIACS}}\\[0.5em]
{\normalsize Based on lectures by Elena Raponi}\\[0.2em]
{\normalsize Academic Year 2025--2026}

\vspace{1.5cm}

\begin{tcolorbox}[colback=boxgray, colframe=leidenblue, width=\textwidth, arc=3mm]
\centering
\textbf{Course Coverage}
\smallskip

\textit{Biological foundations $\cdot$ McCulloch--Pitts \& Perceptron $\cdot$ Hopfield Networks $\cdot$
Multi-layer Perceptron \& Backpropagation $\cdot$ Convolutional Neural Networks $\cdot$
Recurrent Neural Networks \& LSTM $\cdot$ Generative Models $\cdot$ Universal Approximation Theorem}
\end{tcolorbox}
\end{center}

\newpage

% ===========================================================
\section{Introduction and Fundamentals}
% ===========================================================

\subsection{What is Neural Computing?}

Neural Computing studies computational models inspired by the human brain. The discipline focuses on:
\begin{itemize}[leftmargin=*]
    \item Understanding the \textbf{fundamental principles} of neural computation.
    \item Learning the major neural network models (convolutional, recurrent, generative).
    \item Practical implementation using frameworks (e.g., PyTorch).
    \item Tackling computational tasks such as image classification, time-series forecasting, and generation.
\end{itemize}

\subsection{Biological Background}

\subsubsection{The Biological Neuron}

A \emph{neuron} is a cell that transforms and transmits information via electrical and chemical signals. It is:
\begin{itemize}[leftmargin=*]
    \item \textbf{Electrically excitable} (or inhibited).
    \item \textbf{Interconnected} via synapses.
    \item Capable of rapidly propagating signals over large distances.
\end{itemize}

\begin{definitionbox}[title=Neuron Morphology]
\begin{description}
    \item[Dendrites] Receive inputs from other neurons.
    \item[Axon] Carries the neuronal output to other cells.
    \item[Soma] The cell body, containing the cell nucleus.
    \item[Synapses] Connection points between two neurons.
\end{description}
\end{definitionbox}

\subsubsection{Physical Features and Neurotransmission}

\textbf{Ion channels on the membrane} include pumps of ions ($\mathrm{Na}^+, \mathrm{K}^+, \mathrm{Ca}^{2+}, \mathrm{Cl}^-$).

\textbf{Membrane potential} is the difference in electrical potential between the interior of a neuron and the extracellular medium:
\begin{itemize}[leftmargin=*]
    \item \textbf{Depolarization}: $\mathrm{Na}^+$ ions flow inside the cell, making the membrane potential \emph{less negative} $\Rightarrow$ \textbf{excitatory}.
    \item \textbf{Hyperpolarization}: $\mathrm{K}^+$ ions flow out, making the membrane potential \emph{more negative} $\Rightarrow$ \textbf{inhibitory}.
\end{itemize}

\textbf{Action potential (spike)} is the only signal that propagates over long distances and carries information.

\subsubsection{Why the Brain Inspires AI}

\begin{itemize}[leftmargin=*]
    \item Robust and fault tolerant.
    \item Flexible (adjusts to new environments by learning).
    \item Handles fuzzy/noisy/probabilistic information.
    \item Highly parallel.
    \item Compact and energy-efficient.
\end{itemize}

\subsection{Neural Dynamics}

\begin{itemize}[leftmargin=*]
    \item \textbf{Absolute refractory period}: the minimal temporal distance between two spikes.
    \item \textbf{Relative refractory period}: a phase where firing a spike is difficult but not impossible.
    \item \textbf{Postsynaptic Potential (PSP)}: transient impact of incoming spikes; attenuates over time; additive in the model.
    \item A spike is generated if the cumulated membrane potential \emph{reaches the threshold from below}.
\end{itemize}

\subsection{Artificial Neurons}

\subsubsection{The McCulloch--Pitts Model (1943)}

\begin{definitionbox}[title=McCulloch--Pitts Neuron]
\[
O = \sgn\!\left(\sum_i w_i x_i\right), \qquad \sgn(u) =
\begin{cases}
+1 & \text{if } u \geq 0\\
-1 & \text{if } u < 0
\end{cases}
\]
\end{definitionbox}

In this model:
\begin{itemize}[leftmargin=*]
    \item \textbf{Spikes} $\to$ interpreted as spike rates (firing rate).
    \item \textbf{Synaptic strength} $\to$ translated as \emph{weights}.
    \item \textbf{Excitation} = positive product between input and corresponding weight.
    \item \textbf{Inhibition} = negative product between input and corresponding weight.
\end{itemize}

\subsubsection{The Perceptron}

A perceptron is an artificial neuron with a sign (or Heaviside step) activation function, also called a \textbf{threshold unit}:
\[
O = \sgn(\vw^\top \vx + b)
\]
It is a \textbf{linear binary classifier} with decision boundary $\vw^\top \vx + b = 0$.

\textbf{Generalization}: replace the threshold unit with non-linear functions to get real-valued outputs.

\subsection{Network Architectures Overview}

\begin{itemize}[leftmargin=*]
    \item \textbf{Feed-forward Network}: layered, \emph{no cycles}. Topologies: fully connected, randomly connected, optimized (architecture search).
    \item \textbf{Recurrent Network}: contains cycles, is a \emph{dynamical system}, internal state evolves over time; training is more difficult.
    \item \textbf{Hopfield Network}: an associative memory that can store patterns and retrieve them from noisy inputs (John J. Hopfield won the 2024 Nobel Prize in Physics for this contribution).
\end{itemize}

\newpage
% ===========================================================
\section{The Hopfield Network}
% ===========================================================

\subsection{Motivation: Associative Memory}

\begin{definitionbox}[title=Associative Memory]
Also known as \emph{content-addressable memory}, recall is triggered by the \textbf{similarity} between an input pattern and a stored one. It deals with \textbf{incomplete or noisy information}.
\end{definitionbox}

The Hopfield Network (HN) is a \textbf{recurrent neural network} capable of storing and retrieving \emph{multiple} memories. It is inspired by the \textbf{Ising model} from physics, which describes the magnetic behavior of ferromagnetic materials.

\subsection{Structure of a Hopfield Network}

In a Hopfield Network, each unit has three primary characteristics:

\begin{enumerate}[leftmargin=*]
    \item \textbf{Connectivity}: every neuron is connected to all other neurons; each connection has a weight (analogous to synaptic strength) stored in a weight matrix. \textbf{Connections are symmetric}: $w_{ij} = w_{ji}$.
    \item \textbf{Activation}: net input from other neurons through connection weights (analogous to membrane potential).
    \item \textbf{Bipolar state}: output computed via a thresholding function, restricted to values in $\{-1, +1\}$.
\end{enumerate}

\subsection{Formal Definition}

The general dynamic update rule is:
\[
O_i(t+1) = g\!\left(\sum_j w_{ij}\, O_j(t) + b_i\right)
\]
where $g$ can be $\tanh$ or sigmoid. In the discrete/simplified version:
\[
O_i(t+1) = \sgn\!\left(\sum_j w_{ij}\, O_j(t)\right)
\]

\subsection{Stability Condition and Attractors}

\begin{definitionbox}[title=Stable Pattern]
Given a stored pattern $\vec{\xi}$, the stability condition is:
\[
O_i = \sgn\!\left(\sum_j w_{ij}\, \xi_j\right) = \xi_i
\]
\end{definitionbox}

If we set the weights to $w_{ij} \propto \xi_i \xi_j$, the network will reproduce the stored pattern $\vec{\xi}$ (the pattern is an \textbf{attractor}). Stable states partition the input space into \textbf{attracting basins}.

\subsection{Multiple Stored Patterns}

To store multiple patterns $\{\vec{\xi}^1, \vec{\xi}^2, \ldots, \vec{\xi}^P\}$, the weights are defined as a \textbf{superposition of terms}:
\[
w_{ij} = \frac{1}{N}\sum_{\mu=1}^{P} \xi_i^{\mu} \xi_j^{\mu}
\]

This allows:
\begin{itemize}[leftmargin=*]
    \item Associative memory: partial inputs can retrieve full patterns.
    \item Efficient use of the network's capacity.
    \item Pattern recognition applications.
\end{itemize}

\subsection{Storage Capacity}

\begin{keybox}
The \textbf{storage capacity} of a Hopfield Network refers to the maximum number of patterns it can store and accurately retrieve ($P_{\max}$). Empirically, $P_{\max} \approx 0.15\,N$, where $N$ is the number of neurons.
\end{keybox}

If too many patterns are stored, retrieval becomes unreliable due to \textbf{interference} and \textbf{spurious states}.

\textbf{Ways to improve capacity:}
\begin{itemize}[leftmargin=*]
    \item \textbf{Sparse coding}: only a fraction of neurons activate per pattern.
    \item \textbf{Modified learning rules} (e.g., Storkey Learning Rule).
    \item \textbf{More neurons}: increasing $N$ directly increases $P_{\max}$.
    \item \textbf{Hierarchical Networks}: structured memory reduces interference.
\end{itemize}

\subsection{Energy Function}

\begin{theorembox}[title=Energy Function]
\[
E(\vx) = -\frac{1}{2}\vx^\top \mathbf{W}\vx - \vx^\top \vec{b}
\]
\end{theorembox}

\begin{itemize}[leftmargin=*]
    \item \textbf{Stable states} are the minimizers of the energy function.
    \item The Energy function monotonically \textbf{decreases} with each update.
    \item One can use \emph{automatic gradient descent} on this formula to find the weights.
\end{itemize}

\newpage
% ===========================================================
\section{The Simple Perceptron}
% ===========================================================

\subsection{Hebbian Learning}

\begin{keybox}[title=Hebb's Postulate (1949)]
``When an axon in cell $A$ is near enough to \emph{excite cell $B$ and repeatedly and persistently takes part in firing it}, some growth process or metabolic change takes place in one or both cells such that \emph{$A$'s efficiency in firing $B$ is increased}.''
\end{keybox}

In modern terms: \textbf{when two neurons fire together, the connection between them is strengthened}. This is the fundamental operation for learning and memorization.

\[
\Delta w_1(t) \propto x(t)\, y(t)
\]

\subsection{Network Components}

An \textbf{Artificial Neural Network} is characterized by three aspects:
\begin{itemize}[leftmargin=*]
    \item \textbf{Architecture/topology}: nodes and connections with corresponding weights.
    \item \textbf{Node characteristics}: notably, the activation function.
    \item \textbf{Learning rules}: algorithms used to train the network to produce desired outputs for given inputs.
\end{itemize}

\subsection{Activation Functions}

\begin{center}
\begin{tabular}{lll}
\toprule
\textbf{Name} & \textbf{Equation} & \textbf{Range} \\
\midrule
Identity & $f(x) = x$ & $(-\infty,\infty)$ \\
Binary step (Heaviside) & $f(x)= 0$ if $x<0$, $1$ otherwise & $\{0,1\}$ \\
Logistic (Sigmoid) & $\sigma(x) = \frac{1}{1+e^{-x}}$ & $(0,1)$ \\
TanH & $\tanh(x) = \frac{e^x - e^{-x}}{e^x + e^{-x}}$ & $(-1,1)$ \\
ReLU & $f(x) = \max\{0, x\}$ & $[0,\infty)$ \\
GELU & $f(x) = x\,\Phi(x)$ & $(\approx-0.17,\infty)$ \\
SoftPlus & $f(x) = \ln(1+e^x)$ & $(0,\infty)$ \\
ELU & $\alpha(e^x-1)$ for $x\le0$, $x$ otherwise & $(-\alpha, \infty)$ \\
\bottomrule
\end{tabular}
\end{center}

\textbf{Three classical neuron types:}
\begin{itemize}[leftmargin=*]
    \item \textbf{Threshold neuron}: Heaviside / sign function.
    \item \textbf{Saturation neuron}: piecewise-linear function.
    \item \textbf{Sigmoidal neuron}: logistic function or hyperbolic tangent.
\end{itemize}

\subsection{Learning Rule of the Simple Perceptron}

\begin{definitionbox}[title=Learning vs. Training]
\textbf{Learning} = ability to modify proper characteristics (e.g., weights and biases) according to changes in the environment.

\textbf{Training} is typically accomplished via examples; network parameters are adapted using a learning algorithm.

\textbf{Epoch} = a complete run where the entire training set is presented and used by the learning algorithm.
\end{definitionbox}

The simple perceptron has only input and output layers (no hidden layers). For each output $O_i$ ($i=1,\ldots,M$):
\[
O_i = \sgn(\vec{w}^{i\top} \vec{\xi} + b_i)
\]

\textbf{Threshold incorporation}: the bias $b_i$ can be incorporated in the scalar product by introducing an additional input node fixed to 1.

\subsubsection{Hebbian Learning for the Perceptron}

Given a dataset $\{(\vx^p, y^p)\}_{p=1}^{P}$ with $\vx^p\in\R^d$, $y^p\in\{-1,+1\}$:
\[
O^p = \sgn(\vec{w}^\top \vec{x}^p + b)
\]

The goal: find $\vec{w}$ and $b$ such that $O^p = y^p$ for all $p$.

\begin{keybox}[title=Update Rule]
\[
\vec{w}' = \vec{w} + \Delta\vec{w}, \qquad
\Delta w_j = \begin{cases}
2\eta\, y_p\, x_j^p & \text{if } O_p \neq y_p\\
0 & \text{otherwise}
\end{cases}
\]
where $0 < \eta < 1$ is the \textbf{learning rate}.
\end{keybox}

\textbf{Equivalent update formulas:}
\[
\Delta_j = \eta(1 - O_p\, y_p)\, x_j^p\, y_p, \qquad
\Delta_j = \eta(y_p - O_p)\, x_j^p
\]

\subsection{Linear Separability}

\begin{definitionbox}[title=Linear Separability]
Data are \emph{linearly separable} when there exists a weight vector $\vec{w}$ such that:
\[
\vec{w}^\top \vx > 0 \text{ if } y = +1, \qquad \vec{w}^\top \vx < 0 \text{ if } y = -1
\]
\end{definitionbox}

Using $\hat{x}^p = y_p\, x^p$, the condition becomes:
\[
\vec{w}^\top \hat{x}^p > 0, \quad p = 1,\ldots,P
\]

\subsubsection{The AND Function (linearly separable)}

A valid solution: $w_1 = 1$, $w_2 = 1$, $w_3 = \theta = 1.5$.

\subsubsection{The XOR Problem (NOT linearly separable)}

\begin{warningbox}
The XOR function CANNOT be solved by a single perceptron! The constraints lead to a contradiction:
\begin{itemize}[leftmargin=*]
    \item From rows (d) and (b): $w_1 + w_2 < w_3 < -w_1 + w_2$, so $w_1 < 0$.
    \item From rows (a) and (c): $-w_1 - w_2 < w_3 < w_1 - w_2$, so $w_1 > 0$.
\end{itemize}
$\Rightarrow$ contradiction!
\end{warningbox}

\textbf{Solution}: use \emph{two} decision boundaries $g_1(\vx)$ and $g_2(\vx)$, and pass them through threshold functions $y_1 = \sgn(g_1(\vx))$, $y_2 = \sgn(g_2(\vx))$. This maps the input to a new space where the classes are linearly separable. This is the motivation for the \textbf{multi-layer perceptron}.

\subsection{Perceptron Optimality}

Instead of just requiring the sign to be correct, demand a \textbf{margin}:
\[
y_p\, \vec{w}^\top \vx^p > N\rho
\]
for some fixed positive $\rho$ ($N$ keeps the margin fixed for any number of inputs).

The update rule becomes:
\[
\Delta_j = \eta\, \Theta(N\rho - y_p\, \vec{w}^\top \vx^p)\, x_j^p\, y_p
\]
where $\Theta$ is the unit step function.

\begin{keybox}[title=Optimal Weight Vector (Support Vector Machine Idea)]
The optimal $\vec{w}$ \textbf{maximizes the margin} (the distance from the decision boundary to the closest correctly-classified point):
\[
\vec{w}^* = \arg\max_{\vec{w}} M(\vec{w}) = \arg\max_{\vec{w}} \min_{p\in[1..P]} \frac{\vec{w}^\top \hat{x}^p}{\|\vec{w}\|}
\]
This is a \textbf{maximin problem} -- the key idea behind \emph{Support Vector Machines}.
\end{keybox}

\newpage
% ===========================================================
\section{Multi-Layer Perceptron and Backpropagation}
% ===========================================================

\subsection{From Perceptron to MLP}

\begin{itemize}[leftmargin=*]
    \item The \textbf{hidden layer} functions to \textbf{extract features}.
    \item The \textbf{output layer} implements the \textbf{decision function}.
    \item By adding up to 2 hidden layers of perceptrons, one can represent any union of intersections of half-spaces (\textbf{Universal Approximation}).
\end{itemize}

\subsection{Deep Networks vs Shallow ML}

\textbf{Deep Neural Networks} are simply feedforward networks with many hidden layers. The crucial difference from traditional ML lies in \textbf{how} they process information:
\begin{itemize}[leftmargin=*]
    \item Deep networks learn \textbf{hierarchical representations} from data.
    \item They \textbf{automatically extract complex features} at different layers (low-level $\to$ high-level).
\end{itemize}

\subsection{Recap: Derivatives and Gradient}

\textbf{Univariate derivative:}
\[
f'(x) = \frac{df}{dx} = \lim_{\Delta x \to 0}\frac{f(x+\Delta x)-f(x)}{\Delta x}
\]

\textbf{Partial derivative} for multivariate $f:\R^n \to \R$:
\[
\frac{\partial f}{\partial x_i} = \lim_{\varepsilon \to 0}\frac{f(\vx + \varepsilon \vec{e}_i) - f(\vx)}{\varepsilon}
\]

\textbf{Gradient} (collects all partial derivatives):
\[
\nabla f = \left(\frac{\partial f}{\partial x_1}, \ldots, \frac{\partial f}{\partial x_n}\right)^\top
\]
The gradient is the \textbf{steepest ascending direction}.

\subsubsection{Chain Rule}

\textbf{Univariate}: $(g(f(x)))' = g'(f(x))\, f'(x)$

\textbf{Multivariate}: for $O = g(\vec{f}(\vx))$,
\[
\frac{\partial O}{\partial x_i} = \sum_{k=1}^d \frac{\partial O}{\partial f_k}\frac{\partial f_k}{\partial x_i}
\]

\subsection{Gradient Descent}

For a perceptron, define the squared-error loss:
\[
L(\vec{w}, \vx^p) = \tfrac{1}{2}(y_p - \vec{w}^\top \vx^p)^2
\]

The gradient w.r.t. $w_k$ is:
\[
\frac{\partial L}{\partial w_k} = -(y_p - O_p)\, x_k^p
\]

\begin{keybox}[title=Gradient Descent Update]
\[
w_k \leftarrow w_k - \eta\, \frac{\partial L}{\partial w_k}
\quad\Longleftrightarrow\quad
\Delta w_k = \eta(y_p - O_p)\, x_k^p
\]
This re-discovers \textbf{Hebb's rule}!
\end{keybox}

The \textbf{step size $\eta$} is critical:
\begin{itemize}[leftmargin=*]
    \item Too small ($\eta = 0.01$): convergence is slow but stable.
    \item Too large ($\eta = 0.041$): oscillation and divergence.
    \item Moderate ($\eta = 0.035$): zig-zag convergence.
\end{itemize}

\subsection{The Backpropagation Algorithm}

\subsubsection{Notation}

\begin{itemize}[leftmargin=*]
    \item $w_{ij}^{(k)}$: weight for node $i$ in layer $k$ for input $j$.
    \item $b_i^{(k)}$: bias for node $i$ in layer $k$.
    \item $h_i^{(k)}$: sum of inputs for node $i$ in layer $k$.
    \item $O_i^{(k)}$: output for node $i$ in layer $k$.
    \item $n_k$: number of nodes in layer $k$.
\end{itemize}

Recall:
\[
O_i^{(k)} = g(h_i^{(k)}) = g\!\left(\sum_{j=1}^{n_{k-1}} w_{ij}^{(k)}\, O_j^{(k-1)} + b_i^{(k)}\right)
\]

\subsubsection{Step-by-Step Derivation}

\textbf{Step 1 -- Computation Graph}: represent inputs $x$, weights $w$, biases $b$, hidden activations $h$, outputs $O$, label $y$, and loss $L$ as a graph.

\textbf{Step 2 -- Univariate Chain Rule}:
\begin{align*}
\frac{\partial L}{\partial w} &= \frac{\partial L}{\partial O}\frac{\partial O}{\partial h}\frac{\partial h}{\partial w} = \delta\, x \\
\frac{\partial L}{\partial b} &= \delta
\end{align*}
where the \textbf{error term} is $\delta = \frac{\partial L}{\partial O}\frac{\partial O}{\partial h} = \frac{\partial L}{\partial O}\, g'(h)$.

\textbf{Step 3 -- Error term for layer $k$}:
\[
\delta_i^{(k)} = \frac{\partial L}{\partial h_i^{(k)}} = g'(h_i^{(k)})\,\frac{\partial L}{\partial O_i^{(k)}}
\]

The gradients for the weight and bias are:
\[
\frac{\partial L}{\partial w_{ij}^{(k)}} = O_j^{(k-1)}\,\delta_i^{(k)} \quad\text{(Hebb's rule!)}, \qquad
\frac{\partial L}{\partial b_i^{(k)}} = \delta_i^{(k)}
\]

\textbf{Step 4 -- Backpropagating the error}:
\[
\delta_j^{(k-1)} = g'(h_j^{(k-1)})\sum_{i=1}^{n_k} w_{ij}^{(k)}\,\delta_i^{(k)}
\]

\subsubsection{The BP Algorithm Summary}

\begin{tcolorbox}[colback=boxgray, colframe=leidenblue, title=\textbf{Backpropagation Algorithm}]
\textbf{1. Forward phase}: for each input-output pair $(\vx, y)$, starting from the first layer, calculate and store the quantities $O_i^{(k)}$ and $h_i^{(k)}$ for each node.

\smallskip
\textbf{2. Backward phase}: for the output layer, compute:
\[
\delta^{\text{out}} = g'(h^{\text{out}})(y - O)
\]
While moving back to the first layer:
\begin{itemize}[leftmargin=*]
    \item Compute the gradient for each weight and bias:
    $\frac{\partial L}{\partial w_{ij}^{(k)}} = O_j^{(k-1)}\,\delta_i^{(k)}$.
    \item Backpropagate the error: $\delta_j^{(k-1)} = g'(h_j^{(k-1)})\sum_{i=1}^{n_k} w_{ij}^{(k)}\delta_i^{(k)}$.
    \item $k \leftarrow k - 1$.
\end{itemize}

\smallskip
\textbf{3. Gradient descent}: for each weight,
\[
w_{ij}^{(k)} \leftarrow w_{ij}^{(k)} - \eta\,\frac{\partial L}{\partial w_{ij}^{(k)}}
\]
\end{tcolorbox}

\subsection{Variants of Gradient Descent}

\begin{enumerate}[leftmargin=*]
    \item \textbf{Batch Gradient Descent}: use the whole dataset to compute the average gradient.
    \begin{itemize}
        \item \textbf{Pros}: precise, less random.
        \item \textbf{Cons}: additional storage.
    \end{itemize}
    \item \textbf{Stochastic Gradient Descent (SGD)}: take a random data point and compute the gradient with it.
    \begin{itemize}
        \item \textbf{Pros}: less storage, superior for redundant data.
        \item \textbf{Cons}: increased oscillations.
    \end{itemize}
    \item \textbf{Mini-batch Gradient Descent}: take a random subset (batch) and compute the gradient only on this subset.
\end{enumerate}

\subsubsection{Key Terminology}

\begin{itemize}[leftmargin=*]
    \item \textbf{Batch size}: total number of training data points in a single batch (commonly powers of 2: 32, 64, 128).
    \item \textbf{Iteration/step}: the number of batches needed to complete one epoch.
    \item \textbf{Epoch}: one forward+backward pass over the whole training set.
\end{itemize}

\subsection{Overfitting and Generalization}

\begin{itemize}[leftmargin=*]
    \item \textbf{Underfitting}: model is too simple to capture the data.
    \item \textbf{Good fit / Robust}: balances training and validation loss.
    \item \textbf{Overfitting}: model memorizes the training data but fails on validation/test data.
\end{itemize}

\textbf{Early stopping}: monitor validation loss; stop training when it starts to increase.

\subsection{Advanced Optimizers}

Vanilla gradient descent can get stuck in local optima or plateaus. Modern optimizers improve convergence:
\begin{itemize}[leftmargin=*]
    \item \textbf{Momentum}, Nesterov Accelerated Gradient.
    \item \textbf{Adagrad}, Adadelta, RMSprop.
    \item \textbf{Adam}, AdaMax, Nadam, AMSGrad.
\end{itemize}

\subsubsection{Gradient Descent with Momentum}

\begin{keybox}
\[
v_{ij}^{(k)} \leftarrow \gamma\, v_{ij}^{(k)} + \eta\,\frac{\partial L}{\partial w_{ij}^{(k)}}
\]
\[
w_{ij}^{(k)} \leftarrow w_{ij}^{(k)} - v_{ij}^{(k)}
\]
$v_{ij}^{(k)}$ is initialized to zero; $\gamma, \eta \in (0,1)$.
\end{keybox}

Momentum accumulates a velocity vector that smooths the update direction and helps escape shallow local minima.

\newpage
% ===========================================================
\section{Convolutional Neural Networks (CNNs)}
% ===========================================================

\subsection{Motivation: Computer Vision}

\textbf{Sight is one of the most important senses.} Computer Vision tackles:
\begin{itemize}[leftmargin=*]
    \item Image classification.
    \item Object detection.
    \item Segmentation (semantic / instance).
    \item Image captioning.
\end{itemize}

\textbf{Applications}: self-driving cars, medical imaging (cancer detection), accessibility (helping visually-impaired runners).

\subsection{How Computers See Images}

\begin{itemize}[leftmargin=*]
    \item \textbf{Gray-scale images}: a 2D matrix of pixel values (single channel, e.g., $[0, 255]$).
    \item \textbf{RGB images}: 3 channels stacked into a \textbf{tensor}.
\end{itemize}

\subsubsection{Benchmark Datasets}

\begin{itemize}[leftmargin=*]
    \item \textbf{MNIST}: 60,000 training + 10,000 test handwritten digits, $28\times28$ pixels.
    \item \textbf{CIFAR-10}: 50,000 training + 10,000 test images, $32\times32$, 10 classes, 6000 images per class.
\end{itemize}

\subsection{Why Not Use a Fully Connected NN?}

\begin{warningbox}
Naive approach: flatten the image into a vector and feed it to a fully-connected MLP. \textbf{Problems}:
\begin{itemize}[leftmargin=*]
    \item Number of weights ($H \times W$) explodes!
    \item Geometric/spatial information is destroyed.
\end{itemize}
\end{warningbox}

\subsection{Spatial Structure via Convolution}

\begin{keybox}[title=Winning Strategy]
Connect \emph{patches} of input to neurons in the hidden layer:
\begin{itemize}[leftmargin=*]
    \item One node responds only to nodes in a patch of the previous layer.
    \item A single patch maps to a single neuron in the next layer.
    \item Slide the window to define connections.
\end{itemize}
\end{keybox}

\subsubsection{Convolution: Feature Extraction}

\begin{definitionbox}[title=Convolution]
\begin{itemize}[leftmargin=*]
    \item \textbf{Filter (kernel)} of size $H_k \times W_k$: a small matrix of weights.
    \item The \textbf{same filter} is applied to all patches in the input (parameter sharing).
    \item \textbf{Stride}: how many pixels the kernel slides for the next patch (often 1).
    \item Each filter spatially shares the same parameters.
\end{itemize}
\end{definitionbox}

For an input $(H_{\text{in}}, W_{\text{in}})$ and kernel $(H_k, W_k)$:
\[
\text{Output size} = (H_{\text{in}} - H_k + 1, W_{\text{in}} - W_k + 1)
\]

\subsubsection{Local Connectivity}

For each neuron at position $(p, q)$ in the hidden layer:
\[
\sum_{i=1}^{4}\sum_{j=1}^{4} w_{ij}\, x_{i+p, j+q} + b
\]
For a $4 \times 4$ filter.

\subsection{Padding}

\textbf{Problem}: information at the edges may be lost.

\textbf{Solution}: add a border of zeros (padding) of size $(H_p, W_p)$. The output becomes:
\[
(H_{\text{in}} - H_k + H_p + 1, W_{\text{in}} - W_k + W_p + 1)
\]

\subsection{Multiple Channels}

For RGB images:
\begin{itemize}[leftmargin=*]
    \item Use an \emph{individual filter} (of the same size) for each channel.
    \item Sum the resulting pixel values from all channels.
\end{itemize}

\subsection{Multiple Filters}

Multiple filters extract multiple features. Output volume dimensions:
\[
h \times w \times d
\]
where $h, w$ are spatial dimensions and $d$ = number of filters (\textbf{depth}). The receptive field is the location in the input image that a node is path-connected to.

\subsection{Hand-Crafted Filters and Feature Maps}

\textbf{Examples of kernels}:
\begin{itemize}[leftmargin=*]
    \item Identity: $\begin{psmallmatrix}0&0&0\\0&1&0\\0&0&0\end{psmallmatrix}$
    \item Edge detection: $\begin{psmallmatrix}-1&-1&-1\\-1&8&-1\\-1&-1&-1\end{psmallmatrix}$
    \item Sharpen: $\begin{psmallmatrix}0&-1&0\\-1&5&-1\\0&-1&0\end{psmallmatrix}$
    \item Box blur: $\tfrac{1}{9}\begin{psmallmatrix}1&1&1\\1&1&1\\1&1&1\end{psmallmatrix}$
\end{itemize}

\textbf{In CNNs we LEARN the filters} rather than hand-design them.

\subsection{Non-Linearity (ReLU)}

After convolution, apply a Rectified Linear Unit:
\[
g(z) = \max(0, z)
\]
This:
\begin{itemize}[leftmargin=*]
    \item Allows the network to model non-linear data.
    \item Operates pixel-by-pixel: negatives become 0 (no detection), positives are preserved.
\end{itemize}

\subsection{Pooling}

\begin{definitionbox}[title=Pooling]
Pooling \textbf{downsamples} the feature map to:
\begin{itemize}[leftmargin=*]
    \item Reduce computational cost.
    \item Extract shape/smooth features.
    \item Maintain spatial invariance.
\end{itemize}
Common variants: \textbf{Max-pooling} (take max), \textbf{Average-pooling} (take mean), Min-pooling.
\end{definitionbox}

\subsection{CNN Structure Overview}

\begin{tcolorbox}[colback=boxgray, colframe=leidenblue, title=\textbf{CNN for Classification}]
\textbf{Feature Learning} (Convolutional Part):
\begin{enumerate}[leftmargin=*]
    \item \textbf{Convolution}: apply filters to generate feature maps.
    \item \textbf{Non-linearity}: typically ReLU.
    \item \textbf{Pooling}: downsample feature maps.
\end{enumerate}
Repeat several times.

\smallskip
\textbf{Classification} (Fully-Connected Part):
\begin{itemize}[leftmargin=*]
    \item Flatten the final feature maps.
    \item Pass through fully-connected layers.
    \item Apply \textbf{softmax} to express output as probability of class membership.
\end{itemize}
\end{tcolorbox}

\textbf{Softmax}: $\sigma(z)_i = \frac{e^{z_i}}{\sum_{j=1}^K e^{z_j}}$

\subsection{Comparison to Human Vision}

\begin{itemize}[leftmargin=*]
    \item \textbf{V1, V2, V4}: visual neuron areas. Receptive fields increase from V1 to V4. Process simple visual objects (curves, edges).
    \item \textbf{IT (inferior temporal cortex)}: processes complex visual objects (e.g., faces).
    \item Analogously, deep CNNs learn hierarchical representations: early layers detect edges, mid layers detect shapes, late layers detect objects.
\end{itemize}

\subsection{Visualization: Grad-CAM}

\textbf{Gradient-weighted Class Activation Mapping}:
\begin{itemize}[leftmargin=*]
    \item Weighted overlay of all feature maps in a CNN.
    \item Each layer is weighted by the magnitude of its gradient.
    \item Highlights which regions of the image were most influential for a particular classification.
\end{itemize}

\subsection{Classical CNN Architectures}

\subsubsection{LeNet-5 (LeCun et al., 1989)}

The first successful deep learning architecture, originally applied to handwritten zip code recognition.

\subsubsection{AlexNet (Krizhevsky et al., 2012)}

Famous for winning the \textbf{ImageNet Large Scale Visual Recognition Challenge}. Marked the breakthrough of deep learning in computer vision.

\subsubsection{VGG (Simonyan \& Zisserman, 2014)}

Developed by the Visual Geometry Group at Oxford. Uses many $3\times3$ convolutions and stacks them deeply.

\subsubsection{Vanishing Gradient and ResNet}

\begin{warningbox}[title=Vanishing Gradient Problem]
More layers $\to$ more expressive power, BUT gradients at the early layers become tiny, making training very difficult.
\end{warningbox}

\textbf{Solution: Residual Blocks (ResNet, He et al., 2016)}
\begin{itemize}[leftmargin=*]
    \item Introduce a \textbf{skip/residual connection}: $F(\vx) + \vx$.
    \item Gradient flow can backpropagate via the skip connections without being downscaled.
    \item Enables training of very deep networks (50, 100, even 1000+ layers).
\end{itemize}

\textbf{ResNet is the most-used backbone} in many modern models:
\begin{itemize}[leftmargin=*]
    \item RCNN (Regions with CNN features) for object detection.
    \item Image captioning.
    \item Visual Transformers.
    \item Visual autoencoders.
\end{itemize}

\newpage
% ===========================================================
\section{Recurrent Neural Networks (RNNs)}
% ===========================================================

\subsection{Motivation: Time Series Forecasting}

Sequential data exhibits \textbf{temporal dependency}. Each data point depends on previous ones.

\textbf{Can we use a multi-layer NN?} Issues:
\begin{itemize}[leftmargin=*]
    \item How long is the past to use? (Fixed window is restrictive.)
    \item Temporal dependencies among inputs are ignored.
    \item The input-output relationship may vary through time.
\end{itemize}

\subsection{Introducing Cycles}

\begin{definitionbox}[title=Recurrent Neuron]
A neuron with a cycle:
\[
O(t) = g(\vec{w}_1^\top \vec{x}(t) + w_2\, O(t-1) + b)
\]
The output \emph{evolves through time}; given a fixed input, it generates a time series as output.
\end{definitionbox}

\textbf{Applications}:
\begin{itemize}[leftmargin=*]
    \item Memory $\to$ Hopfield Network.
    \item Time-series forecasting and anomaly detection.
    \item Language tasks: machine translation, sentiment analysis.
    \item Generative tasks: Boltzmann machine.
\end{itemize}

\subsection{Computation Graph for RNNs}

\textbf{Unfolding} a simple RNN in time gives a feedforward structure with shared weights across time steps.

\subsection{Vanilla RNN}

\begin{definitionbox}[title=Vanilla RNN with Hidden State]
\begin{align*}
h(t) &= \sigma\!\left(W\begin{bmatrix} x(t) \\ h(t-1) \end{bmatrix}\right) \\
O(t) &= g(h(t))
\end{align*}
The \textbf{hidden state} is the internal state of the RNN cell: it represents a short memory of the history.
\end{definitionbox}

The hidden state is \emph{different from the output}: it carries information forward in time even when the output is ignored.

\subsection{Use Case: Sentiment Classification}

\textbf{Input}: words/sentences (e.g., movie reviews). \textbf{Output}: positive or negative sentiment.

Various readout strategies:
\begin{itemize}[leftmargin=*]
    \item Use only the \textbf{last hidden state} $h_n$ to feed a linear classifier.
    \item \textbf{Ignore} intermediate hidden states.
    \item Use a \textbf{sum} of all hidden states: $h = \text{Sum}(h_1, \ldots, h_n)$.
\end{itemize}

\subsection{Training RNNs: Backpropagation Through Time (BPTT)}

\begin{tcolorbox}[colback=boxgray, colframe=leidenblue, title=\textbf{BPTT}]
\begin{enumerate}[leftmargin=*]
    \item \textbf{Unfold (copy)} the RNN in time during the forward pass.
    \item Treat the unfolded RNN as a \textbf{multi-layer NN}.
    \item Weight updates are computed \textbf{for each copy} in the unfolded network, then averaged.
\end{enumerate}
\end{tcolorbox}

\subsection{Issue: Vanishing or Exploding Gradients}

\begin{warningbox}
\begin{itemize}[leftmargin=*]
    \item If the \textbf{largest singular value of $W$ is small}, the vanishing gradients problem occurs.
    \item If the \textbf{largest singular value of $W$ is large enough}, the gradient explodes.
\end{itemize}
\end{warningbox}

This makes vanilla RNNs unable to learn long-term dependencies.

\subsection{Long Short-Term Memory (LSTM) Networks}

\begin{keybox}
\textbf{LSTMs} (Hochreiter \& Schmidhuber, 1997) are gated RNNs introducing a \textbf{cell state} -- the long-term memory.
\end{keybox}

\subsubsection{Cell State and the Constant Error Carousel (CEC)}

The cell state ensures that gradients don't decay:
\[
\vC(t) = f(t) * \vC(t-1) + \text{Input}(t)
\]

If we ignore the input:
\[
\vC(t) = f(t) * f(t-1) * \cdots * f(t-n+1) * \vC(t-n)
\]
And the gradient:
\[
\nabla \vC(t) = f(t) * f(t-1) * \cdots * f(t-n+1) * \nabla \vC(t-n)
\]

\subsubsection{LSTM Gates}

\textbf{Forget Gate}: decides forgetting coefficient using current input + short history:
\[
f(t) = \sigma\!\left(W_f\, [\vh(t-1), \vec{x}(t)] + b_f\right)
\]
Sigmoid is used to keep $f(t) \in (0, 1)$.

\textbf{Input Gate} (self-gated): scales the current input:
\begin{align*}
\text{Input}(t) &= i(t) * \tilde{C}(t) \\
i(t) &= \sigma\!\left(W_i\,[\vh(t-1), \vx(t)] + b_i\right) \\
\tilde{C}(t) &= \tanh\!\left(W_C\,[\vh(t-1), \vx(t)] + b_C\right)
\end{align*}
$\tanh$ is used to keep $\tilde{C}(t)$ bounded.

\textbf{Cell State Update}:
\[
\vC(t) = f(t) * \vC(t-1) + i(t) * \tilde{C}(t)
\]

\textbf{Output Gate} (self-gated): scales the current cell state:
\begin{align*}
o(t) &= \sigma\!\left(W_o\,[\vh(t-1), \vx(t)] + b_o\right) \\
\vh(t) &= o(t) * \tanh(\vC(t))
\end{align*}

\subsubsection{Why Not ReLU in LSTM?}

\begin{warningbox}
Replacing $\tanh$ by ReLU in the input gate is \emph{not safe}: $i(t) \in (0,1)$ but $\tilde{C}(t)$ computed via ReLU can \textbf{explode}. The same applies to the output gate.
\end{warningbox}

\subsubsection{Caveat: Very Long Sequences}

The forget scaler $f(t) \in (0, 1)$. If $f(t) = 0.99$ and $n = 1000$, then $0.99^{1000} \approx 4.317 \times 10^{-5}$ -- still vanishing!

\textbf{Solution for very long sequences}: \textbf{Attention Mechanism}.

\subsection{Beyond LSTM: Seq2Seq Models}

Sequence-to-sequence is arguably the most classical model for NLP. Used for machine translation, multi-step time-series forecasting.

\subsubsection{Encoder}

\begin{align*}
\vec{x}(t) &= \text{Embedding}(\vec{w}(t)) \\
\vec{h}_E(t) &= \text{RNN}_E(\vec{x}(t), \vec{h}_E(t-1)) \\
\vec{s} &= \vec{h}_E(\text{last})
\end{align*}

\begin{itemize}[leftmargin=*]
    \item Word embedding: typically, a fully-connected NN.
    \item The initial hidden state $\vec{h}(0)$ is usually randomized.
\end{itemize}

\subsubsection{Decoder}

\begin{align*}
\vec{h}(0) &= \vec{h}_E(\text{last}) \\
\vec{h}_D(t) &= \text{RNN}_D(O(t-1), \vec{h}_D(t-1)) \\
O(t) &= \text{Softmax}(\vec{h}_D(t))
\end{align*}

\begin{itemize}[leftmargin=*]
    \item Take $\vec{s} = \vec{h}_E(\text{last})$ to initialize the decoder's hidden state.
    \item Input is the decoded/predicted symbol from the previous step.
\end{itemize}

\newpage
% ===========================================================
\section{Generative Models}
% ===========================================================

\subsection{Introduction}

Formally, generative modeling is \textbf{unsupervised learning}. The goal: create \emph{new forms of data} learned from existing ones, with \textbf{novelty}.

\textbf{Applications}:
\begin{itemize}[leftmargin=*]
    \item Image synthesis.
    \item Image captioning.
    \item Video generation.
    \item Language generation (e.g., GPT models).
    \item Generating reports, slides, etc.
\end{itemize}

\subsection{Autoencoders}

\begin{definitionbox}[title=Autoencoder]
Two (sub-)networks:
\begin{itemize}[leftmargin=*]
    \item \textbf{Encoder} $g_\phi$: obtains the \textbf{latent low-dimensional representation} of the data.
    \item \textbf{Decoder} $f_\theta$: \textbf{reconstructs} the data from the latent code.
\end{itemize}
Ideally $\vx \approx \vx'$ (input $\approx$ reconstruction).
\end{definitionbox}

\textbf{Choice of encoder}:
\begin{itemize}[leftmargin=*]
    \item For images: a CNN.
    \item For time-series: a multi-layer LSTM.
    \item For tabular data: a multi-layer perceptron.
\end{itemize}
$\phi$ represents all trainable parameters (weights and biases).

\textbf{Decoder for image data} requires upsampling. Two options:
\begin{itemize}[leftmargin=*]
    \item \textbf{Upsampling (interpolation)}.
    \item \textbf{Transposed convolution}: spreads each input pixel value over a larger area; can be visualized as inserting zeros between elements and then applying a standard convolution.
\end{itemize}

\subsubsection{Loss Function -- Reconstruction Error}

\[
\mathcal{L}_{\text{AE}}(\phi, \theta) = \frac{1}{n}\sum_{i=1}^n \left[\vx^{(i)} - f_\theta\!\left(g_\phi(\vx^{(i)})\right)\right]^2
\]

Mathematically, we try to learn an identity function: $\text{id} \approx f_\theta \circ g_\phi$.

\subsubsection{Generating New Data}

\begin{enumerate}[leftmargin=*]
    \item After training, sample a random point in the latent space.
    \item Decode this random latent point.
    \item A new image!
\end{enumerate}

\subsection{Variational Autoencoders (VAEs)}

\begin{warningbox}
Standard autoencoders learn a compressed latent representation $z$. However, the latent space may be \textbf{irregular or discontinuous}, and small changes in $z$ may produce unrealistic outputs.
\end{warningbox}

\begin{keybox}[title=VAE Key Idea]
Instead of mapping an input to a single latent point, a VAE learns a \textbf{probability distribution}: $q_\phi(z|\vx)$ parameterized by mean $\mu$ and standard deviation $\sigma$:
\[
z = \mu + \sigma \odot \epsilon, \qquad \epsilon \sim \mathcal{N}(0, I)
\]
\end{keybox}

\textbf{Benefits}:
\begin{itemize}[leftmargin=*]
    \item Smooth latent space.
    \item Meaningful interpolation between samples.
    \item Better generation of new, unseen data.
\end{itemize}

\textbf{Loss function}:
\[
\mathcal{L}_{\text{VAE}} = \underbrace{\mathcal{L}_{\text{reconstruction}}}_{\text{reconstruct input well}} + \underbrace{D_{KL}(q(z|\vx)\,\|\,p(z))}_{\text{keep latent space regular}}
\]
The KL divergence measures how different the encoder's distribution is from a standard normal distribution.

\subsection{Generative Adversarial Networks (GANs)}

Two (sub-)networks compete:
\begin{itemize}[leftmargin=*]
    \item A \textbf{discriminator $D$} (an image classifier): determines if an image is real or synthetic.
    \item A \textbf{generator $G$}: creates synthetic images from noise, trying to make them realistic enough to fool $D$.
\end{itemize}

\textbf{The process}:
\begin{enumerate}[leftmargin=*]
    \item $G$ generates synthetic samples starting from random latent noise vectors.
    \item $D$ receives a mix of real and synthetic images, attempts to distinguish them.
    \item Based on $D$'s success/failure, both networks adjust:
    \begin{itemize}
        \item If $D$ is correct, $G$ improves to create better fakes.
        \item If $D$ is wrong, $D$ learns to avoid similar mistakes.
    \end{itemize}
    \item $D$ is rewarded for correct predictions, $G$ is rewarded for $D$'s errors.
\end{enumerate}

The process continues until both networks reach an \textbf{optimal balance (Nash equilibrium)}.

\subsection{Modern Approaches: Diffusion \& Transformers}

\subsubsection{Diffusion Models}

\begin{itemize}[leftmargin=*]
    \item Modern image generation is dominated by \textbf{diffusion models} much more than GANs.
    \item GANs are historically important, but diffusion models are today's state of the art.
    \item The model learns how to \textbf{iteratively denoise} random noise into meaningful data.
    \item \textbf{Forward diffusion}: gradually add noise to data.
    \item \textbf{Reverse diffusion}: learn to remove noise step by step.
\end{itemize}

\subsubsection{Transformers}

\begin{itemize}[leftmargin=*]
    \item Foundation of modern generative AI: GPT, Gemini, Claude, Llama, etc.
    \item Instead of processing data sequentially (like RNNs), transformers use an \textbf{attention mechanism} to determine which parts of the input are most relevant.
\end{itemize}

\newpage
% ===========================================================
\section{Universal Approximation Theorem (UAT)}
% ===========================================================

\subsection{Theorem Statement}

\begin{theorembox}[title=Universal Approximation Theorem]
Any continuous function $f: [0,1]^n \to [0,1]$ can be approximated with \textbf{arbitrary precision} by a neural network with at least \textbf{one hidden layer} with a finite number of weights.
\end{theorembox}

\textbf{Formal version}: suppose we have a function $f(x)$ to compute within some desired accuracy $\epsilon > 0$. It is guaranteed that, by using enough hidden neurons, we can always find a neural network whose output $g(x)$ satisfies:
\[
|g(x) - f(x)| < \epsilon \quad \text{for all } x.
\]

To this end, the number of hidden nodes, learning iterations, weights, and biases needs to be \emph{adapted individually} for the to-be-approximated function.

\subsection{Visual Proof Sketch}

Consider a function with 1 input and 1 output and a NN with: 2 hidden neurons, 1 output neuron, sigmoid activation $\sigma(wx + b)$.

\subsubsection{Step 1: Generating a Step Function}

By using one neuron with a very large $w$ (e.g., $w = 999$), we generate a \textbf{step function}. The position of the step is:
\[
s = -\frac{b}{w}
\]
proportional to $b$, inversely proportional to $w$. From now on, we use just one parameter $s$ for each hidden neuron.

\subsubsection{Step 2: Weighted Output from Hidden Layer}

For the whole network with hidden neurons making steps at $s_1$ and $s_2$, with weights $w_1$ and $w_2$ to the output:
\[
\text{Output} = w_1\, a_1 + w_2\, a_2
\]
This is a piecewise-constant function with values stepping up at $s_1$ and $s_2$.

\subsubsection{Step 3: Creating a ``Bin'' Function}

By setting $w_2 = -w_1$ (i.e., $h = w_1 = -w_2$), we obtain a single rectangular ``bin'' starting at $s_1$ and ending at $s_2$:
\begin{itemize}[leftmargin=*]
    \item Position controlled by $s_1, s_2$.
    \item Height controlled by $h$.
\end{itemize}

\subsubsection{Step 4: Multiple Bins $\to$ Function Approximation}

By adding pairs of hidden neurons, we create multiple bins:
\begin{itemize}[leftmargin=*]
    \item Divide the interval $[0,1]$ into $N$ sub-intervals.
    \item Use $N$ pairs of hidden neurons to set up peaks at desired heights.
    \item This generates a \textbf{histogram that discretizes the function}.
\end{itemize}

Playing with heights and number of bins, we can approximate any continuous function up to a desired precision. The average deviation between the goal function and the NN approximation can be made arbitrarily small.

\begin{keybox}[title=Practical Implication]
While the UAT guarantees that a \emph{single hidden layer} can approximate any continuous function, in practice we use \textbf{deep architectures} because:
\begin{itemize}[leftmargin=*]
    \item Shallow networks may need exponentially many neurons.
    \item Deep networks learn \textbf{hierarchical features} more efficiently.
    \item Deep networks have better generalization properties.
\end{itemize}
\end{keybox}

\newpage
% ===========================================================
\section{Quick Reference and Exam Tips}
% ===========================================================

\subsection{Key Formulas Cheatsheet}

\subsubsection{Perceptron}
\[
O = \sgn(\vec{w}^\top \vec{x} + b), \qquad
\Delta w_j = \eta(y_p - O_p)\, x_j^p
\]

\subsubsection{Backpropagation}
\[
\delta_i^{(k)} = g'(h_i^{(k)})\,\frac{\partial L}{\partial O_i^{(k)}}, \qquad
\frac{\partial L}{\partial w_{ij}^{(k)}} = O_j^{(k-1)}\,\delta_i^{(k)}
\]
\[
\delta_j^{(k-1)} = g'(h_j^{(k-1)})\sum_{i=1}^{n_k} w_{ij}^{(k)}\delta_i^{(k)}
\]

\subsubsection{Hopfield Network}
\[
w_{ij} = \frac{1}{N}\sum_{\mu=1}^{P} \xi_i^{\mu}\xi_j^{\mu}, \qquad
E(\vec{x}) = -\tfrac{1}{2}\vec{x}^\top \mathbf{W}\vec{x} - \vec{x}^\top\vec{b}
\]

\subsubsection{CNN Output Size}
\[
H_{\text{out}} = H_{\text{in}} - H_k + H_p + 1
\]

\subsubsection{Vanilla RNN}
\[
h(t) = \sigma\!\left(W\begin{bmatrix} x(t) \\ h(t-1) \end{bmatrix}\right), \qquad O(t) = g(h(t))
\]

\subsubsection{LSTM Gates Summary}
\begin{align*}
f(t) &= \sigma(W_f\,[\vh(t-1), \vec{x}(t)] + b_f) && \text{(forget)} \\
i(t) &= \sigma(W_i\,[\vh(t-1), \vec{x}(t)] + b_i) && \text{(input scaler)} \\
\tilde{C}(t) &= \tanh(W_C\,[\vh(t-1), \vec{x}(t)] + b_C) && \text{(candidate)} \\
\vec{C}(t) &= f(t) * \vec{C}(t-1) + i(t) * \tilde{C}(t) && \text{(cell state)} \\
o(t) &= \sigma(W_o\,[\vh(t-1), \vec{x}(t)] + b_o) && \text{(output)} \\
\vh(t) &= o(t) * \tanh(\vec{C}(t)) && \text{(hidden state)}
\end{align*}

\subsubsection{VAE Loss}
\[
\mathcal{L}_{\text{VAE}} = \mathcal{L}_{\text{reconstruction}} + D_{KL}(q(z|\vec{x})\,\|\,p(z))
\]

\subsection{Common Exam Questions}

\textbf{Conceptual:}
\begin{itemize}[leftmargin=*]
    \item Why is XOR not linearly separable? How does an MLP solve it?
    \item What is the difference between a Hopfield Network and a feed-forward NN?
    \item What is the relation between Hebbian learning and gradient descent?
    \item What problem does ResNet solve, and how?
    \item How does an LSTM combat the vanishing gradient problem?
    \item How does a VAE differ from a standard autoencoder?
\end{itemize}

\textbf{Computational:}
\begin{itemize}[leftmargin=*]
    \item Compute the output of a perceptron given weights and inputs.
    \item Compute the output size of a convolutional layer given input, kernel, padding, and stride.
    \item Trace the BPTT for a simple RNN over 2-3 time steps.
    \item Determine whether a given set of patterns is linearly separable.
    \item Calculate the Hopfield weight matrix for a small number of patterns.
\end{itemize}

\textbf{Architecture:}
\begin{itemize}[leftmargin=*]
    \item Sketch the structure of an MLP, CNN, RNN, LSTM, autoencoder, VAE, GAN.
    \item Identify the components of a CNN (conv, pool, FC, softmax).
    \item Identify the LSTM gates and their purposes.
\end{itemize}

\subsection{Important Concepts Summary}

\begin{center}
\begin{tabular}{p{4cm} p{10cm}}
\toprule
\textbf{Concept} & \textbf{Key takeaway} \\
\midrule
McCulloch-Pitts neuron & First computational model of a neuron (1943). \\
Perceptron & Linear binary classifier; cannot solve XOR. \\
Hebbian learning & ``Neurons that fire together, wire together.'' \\
Hopfield Network & Recurrent network for associative memory; capacity $\approx 0.15N$. \\
Backpropagation & Chain rule applied through computation graph to compute gradients. \\
Gradient Descent & Iterative minimization of loss via $w \leftarrow w - \eta\nabla L$. \\
Mini-batch SGD & Trade-off between full-batch and stochastic gradient descent. \\
Overfitting & High training accuracy, low validation accuracy; use early stopping. \\
CNN & Local connectivity + parameter sharing for spatial data. \\
Pooling & Downsampling for invariance and efficiency. \\
ResNet & Skip connections to enable very deep networks. \\
RNN & Cycles in network $\to$ memory of past inputs via hidden state. \\
BPTT & Unfolding the RNN in time and applying BP. \\
LSTM & Gated RNN with cell state to combat vanishing gradients. \\
Seq2Seq & Encoder-decoder for variable-length sequences (translation). \\
Attention & Replace recurrence; foundation of Transformers. \\
Autoencoder & Encoder $\to$ bottleneck $\to$ decoder; learns identity. \\
VAE & Probabilistic latent space; KL regularization. \\
GAN & Generator vs Discriminator adversarial game. \\
Diffusion Model & Learns to iteratively denoise into data; state of the art. \\
UAT & One hidden layer can approximate any continuous function. \\
\bottomrule
\end{tabular}
\end{center}

\vfill
\begin{center}
\rule{0.5\textwidth}{0.4pt}\\
\textit{End of Study Guide -- Good luck with your exam!}\\
\rule{0.5\textwidth}{0.4pt}
\end{center}

\end{document}
